\section{Geometry, compactness, and the martingale observable}
\label{geo:section}
\label{sec:geometry}

Throughout this section \(m>0\), and the physical diameters of the
approximating domains are bounded by a fixed number \(L\).
The restriction on the diameters will be removed in
Section~\ref{loc:section}. Constants may depend on \(m,L\), and the common
angle bound, but never on the lattice, the mesh, or an explored slit.
An arm always means a path of actual sampled open edges of the indicated
color. Identifications used to specify a boundary condition are not edges
of an arm.

\subsection{The two established inputs}

\begin{proposition}[Observable convergence and rectangle crossings]
\label{prop:inputs}
Consider the weights of Section~\ref{mod:section}, with
\(r_\delta/\delta\to m/2\), on isoradial meshes whose rhombus half-angles
belong to \([\eta,\pi/2-\eta]\).
\begin{enumerate}
\item
For marked Carath\'eodory convergence to a bounded simply connected
two-point domain, the reconstructed two-point FK fermion converges
locally uniformly, up to its common sign, to a smooth massive-holomorphic
function \(f\):
\[
 \bar\partial f+\ii m\overline f=0.
\]
Its square integral \(h\), normalized by its two boundary values, is the
unique bounded solution
\[
 \lap h=-4m|\grad h|,\qquad
 h=1\ \hbox{on the wired bank},\qquad
 h=0\ \hbox{on the dual-wired bank}.
\]
The boundary values are prime-end limits away from the two marked ends,
and \(0\le h\le1\). The convergence includes uniform interior bounds
on protected compact subsets.
\item
For every fixed rectangle aspect ratio and every finite physical scale
bound, the crossing probabilities of both colors are bounded away from
zero, uniformly over all boundary conditions and all nondegenerate
lattice rectangles below that scale whose side lengths exceed a fixed
multiple of the mesh. The scale constants are independent of the mesh.
\end{enumerate}
\end{proposition}

The observable convergence, its two-point characterization, and the
rectangle estimates are the parts of \cite[Theorems 1.1 and 1.3,
Definition 5.2, Proposition 5.8]{Park} that we use. Uniqueness in the
full class of bounded scalar solutions, as stated above, follows from
the comparison argument of Proposition~\ref{mod:bvp-unique}.
The abstract angle in
that paper is the angle defined here by the elliptic integral, and its
nome is \(r_\delta\), not the cluster weight. The allowance
\(r_\delta/\delta\to m/2\), rather than an exact equality for every mesh,
is part of its scaling setup. Rescaling a rectangle changes its effective
mass by its physical scale; hence a fixed upper scale bound gives the
uniform version of the second assertion. No convergence theorem for
massive interfaces is included in Proposition~\ref{prop:inputs}.

For later use, we specify the probabilistic normalization of the first
input. Let \(\nu_\xi\) be the unit vector from the primal to
the dual vertex of a corner \(\xi\). For an evaluated corner
\(\xi\), let \(\operatorname{wind}(\gamma^\delta:b^\delta
\rightsquigarrow\xi)\) denote the total tangent turning along the
interface traversed backwards from its terminal corner to \(\xi\).
We fix the terminal phase by the corner projection convention in
\cite[Proposition~2.2 and Lemma~2.5]{Park}. The corner fermion is
\begin{equation}\label{geo:corner-observable}
 F^\delta(\xi)=
 \left(\frac{1}{2\delta\ii\nu_{b^\delta}}\right)^{1/2}
 \E\!\left[
 \mathbf 1_{\{\xi\in\gamma^\delta\}}\,
 \exp\!\left(-\frac{\ii}{2}\,
 \operatorname{wind}(\gamma^\delta:b^\delta\rightsquigarrow\xi)\right)
 \right].
\end{equation}
Indeed, the tangent directions give
\(\exp(\ii\operatorname{wind}(b^\delta\rightsquigarrow\xi))
=\nu_\xi/\nu_{b^\delta}\), so
\(F^\delta(\xi)\in(\ii\nu_\xi)^{-1/2}\R\), as required by
the projection rule. Its square primitive satisfies
\(H(u)-H(w)=2\delta|F^\delta(\langle uw\rangle)|^2\);
at the terminal corner this gives the unit boundary jump.
This also agrees with the normalization in
\cite[Section~2.1.1]{CS}, where \(\nu_{b^\delta}=-\ii\), after
rotation to arbitrary terminal direction.\footnote{For the
primal-to-dual direction \(\nu\) and ordinary tangent turning used here,
the factor \(\ii\) in the terminal phase is forced by the projection
lines. The printed prefactor in \cite[Equation~(2.1)]{Park} omits
this factor; we use its projection and square-integral conventions
to specify the normalization.}
We use one fixed choice of the terminal square root. Changing that choice
multiplies the field by a common sign and has no effect on its square.
Edge reconstruction is a fixed real-linear local projection
of corner values. In particular, while a test edge is still in the
active unexplored domain, its reconstructed value is a martingale for
exploration from the initial corner. The phase is continued from the
original terminal end; it is not reset at the moving tip.

The other established input is the curve theorem of \cite{KS}.
We use its conditional unforced-annulus condition G2, the equivalent
power condition G3, their conformal transport
\cite[Proposition 2.6 and Remark 2.8]{KS}, and joint compactness of curves
and driving functions \cite[Theorem 1.5 and Corollary 1.7]{KS}.
In its annulus formulation the inner circle must meet the current
boundary. A component is unforced when removing it does not separate
the current initial mark from the terminal mark. The condition bounds
the event of crossing in the union of all such components. Its
consequences include a continuous driver, strict capacity growth,
convergence in the curve metric, and attainment of the target only at
the end. We establish the condition for the present model below,
including its graph-boundary and submesh conventions.

\subsection{Exploration domains and a simultaneous arm comparison}

We record a graph convention used throughout the paper. Boundary edges
of the wired and dual-wired banks are regarded as already open in their
respective colors. A turn of the interface exposes the adjacent
diagonal and extends one of those banks. Clip the revealed half-rhombus
along that diagonal, and retain the component containing the continuing
interface. If two boundary occurrences occupy the same physical vertex,
they are first recorded as distinct incident sectors. A forced turn is
traversed without a new random decision.

\begin{lemma}[Exploration and boundary occurrences]
\label{geo:domain-markov}
After any finite sequence of these revelations, the unqueried
continuing problem is an isoradial Dobrushin problem with the same
interior weights. Its two banks are the old banks extended by the
revealed diagonals. All discrepancies between the geometric slit and
this graph domain occur in a bounded number of cell layers along the
slit, together with components cut off from the continuation.
The conditional probabilities of the two outcomes of a genuine new
pairing are uniformly bounded below.
\end{lemma}

\begin{proof}
Pair the medial half-edges locally at each revealed rhombus. Pairings
already fixed can be removed from the product defining the law.
A closed medial loop contributes its fixed factor \(\sqrt2\);
the remaining pairings give the same loop law in the continuing
component. In the cluster formulation this is precisely the spatial
domain Markov property with cluster weight \(2\). The two boundary
colors meet at the current and terminal corners.

Clipping along an open primal diagonal leaves a primal wired edge on
the retained side; clipping along a dual-open diagonal has the dual
interpretation. These are the half-rhombus boundary conventions in the
observable problem of Proposition~\ref{prop:inputs}. If a bank contacts
itself, its different occurrences are read in boundary order. Components
behind the contact have no continuing change point and may be removed.
Thus no new interior weight or extra boundary-change point is introduced.
This proves the assertion even when the physical polygon has pinches:
the simply connected polygon is the one with its incident sectors kept
separate. Equivalently one may draw a thin polygonal collar within the
corresponding half-rhombi. A fixed number of neighboring cell layers
contains the entire change.

Finally, the ratio of the two local edge weights stays in a compact
subset of \((0,\infty)\), by the angle bound and the convergence of
the abstract angles to the geometric angles. Opening an edge changes
the number of FK clusters by either zero or one. The two conditional
odds therefore differ from the weight ratio by at most the fixed
cluster factor \(2\). Both outcomes have a common positive lower bound.
\end{proof}

\begin{lemma}[Comparison for a union of free-bank regions]
\label{geo:arm}
Let \(A=A(z;r,R)\) be a physical annulus with \(R\le L_1\), where \(L_1\)
is fixed. Inside \(A\), let \(\mathcal V\) be a union of regions whose
accessible noncircular banks all have one color, and let \(c\) be the
other color. Circular ends may carry arbitrary induced boundary
identifications. After buffering the two ends by a fixed number of
cell layers, the probability of an actual \(c\)-arm crossing \(A\)
in any of the regions is at most
\begin{equation}
 K\left(\frac{r\vee\delta}{R}\right)^\alpha.
 \label{geo:arm-bound}
\end{equation}
The constants are independent of the number of regions and remain
valid after arbitrary conditioning outside buffered shells. For
specified disjoint buffered annuli the corresponding upper bounds
can be multiplied.
\end{lemma}

\begin{proof}
We first describe the comparison, since simply summing over components
would not prove the assertion. Work in the graph of color \(c\), keeping
boundary-sector copies until the comparison is made. Let \(H\) be the
union of all its portions in the regions under consideration. An
actual \(c\)-path cannot exit \(H\) through a noncircular bank: that bank
is wired in the other color. Consequently every exterior connection
that can influence \(H\) enters through its radial-end vertices.
Conditioning on all edges off the shell may impose any partition of
these vertices. Dominate it by the single partition wiring all radial
ends together, including the inner and outer ends.

Identify copies that project to the same ambient lattice vertex.
For \(q=2\) this is an increase in boundary identification, hence it
increases open-edge connectivity. Embed the resulting graph into the
ambient lattice annulus. Edges missing because of a bank or a deleted
sector may be viewed as conditioned \(c\)-closed. Release those
conditions and add the remaining ambient vertices and edges.
Random-cluster monotonicity couples these operations so that every
open edge of \(H\) remains open in the larger graph. Each counted arm
projects to an actual arm of the ambient graph. The construction
couples the entire union at once and therefore pays no multiplicity
factor.

This also explains why winding around the annulus causes no additional
problem. A winding region changes the partition induced at the radial
ends, but that partition was already dominated by the maximal one.
An artificial wire is used only in the measure and is not counted as
an edge in the arm event. Thus an opposite-color circuit in the
interior of a buffered subannulus blocks every counted arm, even if
the boundary condition artificially connects the two radial ends.

In a subannulus of fixed ratio, cover a circuit by a fixed number of
overlapping rectangles. Proposition~\ref{prop:inputs} bounds below
the conditional probability of each opposite-color crossing under
the worst boundary condition. FKG glues those crossings into a
circuit, giving a probability at least \(c_*>0\), uniformly under
conditioning off that subannulus. Use a family of disjoint buffered
subannuli of fixed ratio between scales \(K_0(r\vee\delta)\) and
\(R/K_0\). Exposing them successively gives an upper bound
\((1-c_*)^N\), with
\(N\ge c\log(R/(r\vee\delta))-C\). This is
\eqref{geo:arm-bound}. The same exposure argument proves the product
claim for disjoint annuli. Restriction to a rectangle inside any
shell induces a boundary partition that is dominated by its fully
wired partition, so arbitrary off-shell conditioning is harmless.
\end{proof}

\begin{remark}
\label{geo:actual-arm}
The noncircular-bank hypothesis is a statement about the accessible
side of the bank. A vertex identified with another occurrence on the
far side of a slit does not supply an accessible wire of color \(c\).
One first works in incident sectors and then makes the monotone
identification in the proof. Likewise an incident determined diagonal
of color \(c\) would expose its own color on an accessible bank; it
cannot be silently placed on a bank declared to have the other color.
Buffering by more than the diameter of the local cell templates removes
the only ambiguity at the circular ends.
\end{remark}

\begin{proposition}[Conditional unforced crossings]
\label{geo:g2}
The regularly drawn medial interfaces in domains of diameter at most
\(L\) satisfy condition G2 of \cite{KS}, uniformly over the allowed
meshes, domains, and thermal parameters.
\end{proposition}

\begin{proof}
At a stopping time, complete the currently exposed tile if the drawing
has been stopped in its interior. This reveals only a bounded number
of local turns. Use the graph domain of
Lemma~\ref{geo:domain-markov}. Consider an annulus whose inner circle
meets its boundary, and a component of its intersection with the domain
that does not separate the two current marks.

Such a component cannot have accessible banks of both colors. Indeed,
an interior path joining the two banks, together with short terminal
pieces in their incident sectors, would be a crosscut joining the two
Dobrushin arcs. The two change points alternate with the ends of that
crosscut on the prime-end circle, so it separates them. Removing the
whole annular component would then also separate them, a contradiction.
An interface crossing in a component with one bank color has an actual
arm of the other color beside its strand. The neighboring diagonals
join in boundary order along the strand; trimming a fixed number of
cells at each circular end removes any terminal ambiguity. Components
with no accessible bank can be assigned either color, since the strand
has both adjacent cluster arms.

Apply Lemma~\ref{geo:arm} separately to the two possible missing colors.
For radii larger than a fixed multiple of the mesh this gives
\(K(r/R)^\alpha\) for the probability of any unforced crossing, with
no dependence on the number of components. For \(R\) exceeding a fixed
multiple of \(L\) a crossing is impossible; thus all rectangle scales
used by the comparison are uniformly bounded.

It remains to give the drawing convention at submesh scales. In each
rhombus use simple polygonal tracks through its side midpoints for its
two possible pairings. Choose the tracks and the boundary half-cell
tracks with a clearance \(c_\eta\delta\) between nonincident pieces.
The initial and terminal tracks leave their boundary corners
transversely. Compactness of the angle interval gives uniform
nondegeneracy of these finitely many types of local templates,
including their lengths, turning angles, and transverse clearances.
This is a perturbation inside a bounded number of cells.

If \(r\ll\delta\), an inner circle meeting the current boundary can be
visited by the future track only within \(O(r)\), measured along the
incident template, of the current or terminal mark. Away from
consecutive pieces this follows from clearance. Along consecutive
pieces it follows from their uniformly nondegenerate angles. In that
local picture both incident banks occur. A subannulus with inner radius
larger than a fixed multiple of \(r\), outer radius smaller than
\(c_\eta\delta\), and still inside the original annulus separates the
mark from the distant continuation. Its component is therefore forced.
Consequently, if \(R/\delta\) is bounded and \(R/r\) is sufficiently
large, there is no unforced crossing at all. If \(R/\delta\) is larger
than a prescribed large constant, use instead
Lemma~\ref{geo:arm} between scale \(K\delta\) and \(R/K\).
Choosing that constant first makes its upper bound smaller than the
constant required in G2. These two cases cover all submesh inner radii.
The argument is conditional on the entire revealed past, as required.
\end{proof}

\subsection{Conformal control of cell changes}

The following elementary use of conformal geometry will always concern
connected \emph{interior} sets and their specified endpoint occurrences.
The distinction matters in a slit domain: the two prime-end copies of
one physical slit point are not a connected set in the prime-end
compactification.

\begin{lemma}[Small connected sets in normalized coordinates]
\label{geo:crosscut}
Let \(V_n\) be simply connected planar domains containing \(z_0\), and
let \(\chi_n:V_n\to\D\) satisfy
\(\chi_n(z_0)=0\) and the same angular normalization. Suppose
\[
 0<c_0\le \crad_{V_n}(z_0)\le C_0<\infty.
\]
If \(E_n\subset V_n\) is connected and
\(\diam E_n\to0\), then \(\diam\chi_n(E_n)\to0\).
The same conclusion holds after adjoining accessible prime-end
endpoints approached along \(E_n\).
In particular, bounded-cell changes of a mesh tending to zero have
vanishing diameter in the normalized disc, provided their incident
sectors are respected.
\end{lemma}

\begin{proof}
We give a length--area proof; in fact only the lower conformal-radius
bound is needed. First let \(V\) be a finite polygon, put \(c=c_0\),
and let \(E\subset V\) be connected with diameter \(d<c/64\).
Koebe's theorem gives \(B(z_0,c/4)\subset V\).
If \(E\) meets \(B(z_0,c/8)\), Cauchy's estimate for the bounded
analytic function \(\chi\), on a slightly larger interior disc,
gives \(\diam\chi(E)\le16d/c\).

Otherwise choose \(x\in E\), put \(r=2d\) and \(R=c/16\).
Then \(E\subset B(x,r)\), whereas \(z_0\notin B(x,R)\).
For a regular radius \(s\in(r,R)\), let
\[
 L(s)=\int_{V\cap\partial B(x,s)}|\chi'(z)|\,|\dd z|
\]
be the \emph{total} image length of all circular components.
Cauchy--Schwarz and polar integration give
\[
 \int_r^R L(s)^2\frac{\dd s}{s}
 \le 2\pi\int_{V\cap\{r<|z-x|<R\}}|\chi'(z)|^2\,\dd A(z)
 \le2\pi^2.
\]
There is therefore a regular \(s\) for which
\[
 L(s)\le \ell:=
 \left(\frac{2\pi^2}{\log(R/r)}\right)^{1/2}.
\]
The source circle is either wholly inside \(V\), or cuts it into
finitely many disjoint crosscuts. Their images have total length
at most \(\ell\). When \(\ell<1/4\), each image crosscut lies within
distance \(\ell\) of one of its circle endpoints. The shorter
boundary-circle arc joining its endpoints lies in that same ball.
The enclosed Jordan region has diameter at most \(2\ell\) and
does not contain \(0\); it is exactly the noncentral shadow of
that crosscut. A full image circle of length \(\ell\) has bounded
interior of diameter at most \(\ell/2\).

The connected set \(\chi(E)\) lies in a component separated from
\(0\) by this finite collection of image crosscuts. Such a
component is contained in the noncentral shadow of one crosscut:
adding disjoint crosscuts one at a time gives a tree of
components, and the first edge on the route from the central
component identifies a separating shadow. Thus
\(\diam\chi(E)\le2\ell\).
This proves a modulus of order
\((\log(c/(32d)))^{-1/2}\), independent of physical area,
diameter, or the number of boundary pieces.

For a general simply connected \(V\), first connect any two
points of \(E\) by a compact path in an arbitrarily small open
neighborhood of \(E\) contained in \(V\). Such a neighborhood
can be chosen connected, as a union of small discs centered on
the connected set \(E\). Its physical diameter is at most
\(d+o(1)\). Exhaust \(V\) by finite Jordan polygons containing
this path and an interior reference disc. Their conformal radii
eventually exceed \(c/2\), and their normalized maps converge
on compact sets. The polygon estimate passes to the two chosen
points. Taking their supremum proves the assertion for \(E\).

Accessible endpoints are handled by truncating their actual
approaching arcs and then taking prime-end limits. For a
boundary-following connected lift in one of the finite polygon
domains used here, write \(f=\chi^{-1}\) and approximate its
parameter image \(C\subset\overline\D\) by \(tC\), \(t\uparrow1\).
The continuous physical boundary extension of \(f\) makes
\(f(tC)\) an interior approximation with the same limiting
physical diameter. This extension need not be injective at a
slit, which is why the lift, not an unlabelled physical boundary
set, was specified. A retained cell sector and a fixed number
of neighboring sectors have physical diameter \(O(\delta)\);
the last assertion follows.
\end{proof}

\begin{proposition}[Joint compactness]
\label{geo:compactness}
Under the common diameter bound, the laws of the interfaces transferred
by their normalized maps to the closed disc are tight in the oriented
curve metric. Every subsequential limit is a chordal Loewner trace with
a continuous driving function and strict capacity growth. It reaches
the target only at the end and is continuously completed in disc
coordinates. Along a further coupling, the curves and their
half-plane drivers converge jointly, with uniform convergence of the
drivers and capacity-parametrized curves on every finite capacity
interval.
\end{proposition}

\begin{proof}
Apply \cite[Theorem 1.5 and Corollary 1.7]{KS} to
Proposition~\ref{geo:g2}. Their modified simple medial drawing agrees
with the separated drawing used above. The changes of drawing have
vanishing disc diameter by Lemma~\ref{geo:crosscut}, and so do not
change the limiting curve law.

The disc maps used in that theorem fix the two endpoints, whereas our
normalization fixes an interior point and angle. Marked Carath\'eodory
convergence makes the intervening disc automorphisms converge to a
nondegenerate automorphism. Thus the two normalizations give the same
tightness and convergence statements. Choose the half-plane maps
\(\phi_\delta\) taking the marks to \(0,\infty\) so that
\(\phi_\delta\to\phi\) on compact subsets; the remaining positive
dilation is fixed by the interior normalization. The joint
half-plane conclusion then has precisely the capacity convention
\(\partial_tg_t=2/(g_t-W_t)\) used here.
\end{proof}

\begin{lemma}[Active graph domains and moving prime ends]
\label{geo:active}
In a coupling as in Proposition~\ref{geo:compactness}, let
\(t_n\to t<\infty\). The active graph domains at the corresponding
exploration times converge in the marked Carath\'eodory sense to
\(O_t\), locally from any protected point in \(O_t\). Their moving
corner converges to the prime end represented by \(W_t\).
The assertion includes contacts of the limiting trace with itself or
with an original bank.
\end{lemma}

\begin{proof}
Work first in the normalized initial disc. The kernel of the geometric
slit domains on the nonswallowed side is identified by joint
Loewner convergence. To compare a graph domain with its geometric
slit, take a compact interior path and a compact neighborhood of it
in \(O_t\). They have positive distance from the limiting past.
For all large \(n\), neither the discrete past nor its cell collar
meets this neighborhood. Join the path to a fixed interior point
near the terminal mark. A sufficiently small terminal neighborhood
is untouched by the past through the given finite time, because
the target is reached only at the end. The remaining connection to
the terminal boundary occurrence can therefore be made in that
unchanged neighborhood. The compact path belongs to the active
graph component.

This proves that every protected compact subset of \(O_t\) is retained.
Conversely, no kernel from such a basepoint can cross the limiting
slit. A component joined only by a neck collapsing to that slit is
not part of its kernel. Thus removal of discarded components and
the cell collar changes do not alter the active kernel.

Kernel convergence alone does not identify the moving prime end.
Choose times \(s>t\), arbitrarily close to \(t\), for which
\(\gamma(s)\) is interior to \(O_t\). Such times exist: a whole
interval of trace confined to the old boundary and hull would add
no hull and hence no capacity, contrary to strict capacity growth.
For fixed \(s\), a compact corridor from \(\gamma(s)\) to an
untouched target neighborhood is retained in the graph domain by
the preceding paragraph. The corresponding discrete continuation
from its moving corner to a point converging to \(\gamma(s)\) is
an actual connected chain in that domain, up to its incident
boundary sectors. Include all forced steps before the next
decision; they remain in the same chain.

By uniform curve convergence, the diameter of this chain in the
initial disc can be made arbitrarily small by first taking \(s-t\)
small and then taking \(n\) large. Normalize the current domains
at a fixed protected bulk point. Their conformal radii there are
uniformly nondegenerate. Lemma~\ref{geo:crosscut} therefore makes
the current-chart diameter of the chain small. Its interior
endpoint converges by kernel convergence, whereas its other end
is the current corner occurrence. As \(s\downarrow t\), the image
of the interior endpoint tends to \(W_t\), by the trace property.
The moving corner consequently tends to the same prime end.
There is no use of marked convergence in applying
Lemma~\ref{geo:crosscut}; it is a uniform statement for connected
interior sets. In particular, near a self-contact the actual
continuation chain selects the correct incident sector.

Finally, returning to physical coordinates uses the compact
convergence of the initial uniformizations. The proof can be
localized at any protected bulk basepoint in \(O_t\), so it does not
require the original basepoint to remain unswallowed.
\end{proof}

\Needspace{13\baselineskip}
\begin{corollary}[Conditional geometric estimates in the limit]
\label{geo:restart}
Every limit of Proposition~\ref{geo:compactness} inherits the conditional
crossing estimates needed below in its natural driving filtration.
They hold after a live stopping time and after localization on a
finite interval on which the relevant Loewner coordinates remain
nondegenerate. In the following two estimates, reset capacity time
to zero at the chosen live stopping time.
\begin{enumerate}
\item In the current half-plane chart, put the tip at \(0\), rescale a
given live point to modulus \(1\), and write \(x\) for its real part.
If its imaginary part is small, so \(x\) is near \(1\) or \(-1\), then
\[
 \Prob\{\text{the future trace visits }B(x,d)\mid\text{past}\}
 \le Kd^\alpha
\]
for all sufficiently small \(d\).
\item For a live physical point \(w\),
\[
 \Prob\left\{
   \inf_{s\text{ before swallowing}}
   \crad_{O_s}(w)
      <\varepsilon\,\crad_{O_0}(w)
   \,\middle|\,\text{past at the restart}\right\}
 \le K\varepsilon^\alpha .
\]
\end{enumerate}
The assertions apply in particular at inverse elapsed-clock times
and at first angular-level crossings used in Section~\ref{eng:section}.
No semimartingale property of the driver is assumed here.
\end{corollary}

\begin{proof}
At a discrete restart, the future graph problem is Dobrushin by
Lemma~\ref{geo:domain-markov}. Proposition~\ref{geo:g2}, conformal
transport, and the equivalence of G2 and G3 give the same conditional
power constants in its current chart. For the first event, use a
half-annulus centered at \(x\), of inner radius \(d\) and a fixed small
outer radius avoiding \(0\). It is unforced for the marks \(0,\infty\),
and its inner circle meets the real boundary. Entry into the ball
requires its crossing. A ball later swallowed without being visited
does not create an entry event.

For the second event, write \(r_0=\crad_{O_0}(w)\).
Koebe's theorem puts the old boundary at distance at least \(r_0/4\).
A decrease of the active conformal radius below \(\varepsilon r_0\)
therefore requires the new trace to approach \(w\) within
\(O(\varepsilon r_0)\). The annulus from that scale to \(r_0/8\) lies
in the old domain. The interface has an actual cluster arm across
it. The blocking-circuit proof of Lemma~\ref{geo:arm}, now in a bulk
annulus, gives the bound. Only physical radii bounded by the original
diameter are used. The infimum stops at swallowing; assigning radius
zero after swallowing would give a different and false assertion.

We explain the limiting conditional statement. First stop before
a protected point has small chart height, before chart coordinates
or the driver leave fixed bounds, and before a fixed capacity
horizon. Joint convergence and Lemma~\ref{geo:active} make the
current charts converge on all compact sets used in the estimate.
Apply the discrete conditional inequality, multiplied by any
bounded nonnegative continuous function of the stopped driver
history. Pass to the limit with slightly enlarged balls and relaxed
thresholds; then decrease those enlargements. A monotone-class
argument gives the conditional inequality for the limiting history.

On such a protected interval, a clock with rate
\(|g_t(\phi(w))-W_t|^{-2}\) is continuous and strictly increasing.
Its inverse stopping times are obtained by convergence of these
clocks. First-crossing times of angular levels can either be taken
at generic levels or approximated on clock-time grids. Continuity,
followed by exhaustion of the protected intervals, gives all
stopping times used here. This procedure concerns only the driver
filtration. It does not presume that conformally mapping a
nonsimple future trace through an old swallowed pocket is a valid
operation.
\end{proof}

\subsection{The signed martingale field}

\begin{proposition}[Continuum fermion martingales]
\label{geo:martingale}
Every subsequential curve law of Proposition~\ref{geo:compactness}
has a joint realization with a complex field \(M_t(w)\), \(w\in O_t\),
such that
\[
 M_t(w)^2=D h_t(w),\qquad
 \bar\partial M_t+\ii m\overline{M_t}=0.
\]
The field is continuous locally in space and time before swallowing,
and is smooth in its spatial variable. On every protected compact
subset of the live domain, after localization, \(M_t(w)\) is a
continuous martingale in the joint natural filtration of the curve
and field. In particular it is a continuous local martingale at
each fixed point before its swallowing time.
\end{proposition}

\begin{proof}
First consider deterministic sequences of active graph domains along
a joint curve--driver convergence as above. Lemma~\ref{geo:active}
and Proposition~\ref{prop:inputs} identify the locally uniform limit
of the squares of their fermions as \(D h_t\). The same reasoning
for every sequence of times \(t_n\to t\) gives convergence locally
along the evolution.

We also need continuity of the limiting square under this domain
variation. It follows directly from the deterministic observable
theorem: approximate each marked continuum domain by an admissible
lattice domain and choose the meshes diagonally. Interior compactness
and uniqueness of the bounded two-point solution identify every
subsequential limit. Equivalently, bounded solutions of the displayed
elliptic equation are compact on interior discs and their prime-end
boundary values are controlled by the same two-point boundary
estimates. Thus the limiting square is continuous in the marked
kernel topology, locally on protected compacts. The assertion applies
uniformly along a compact family of convergent curve--driver paths:
otherwise a contradicting sequence of times and paths would give a
deterministic violating sequence of domains.

It remains to lift convergence of squares to signed martingales.
We give the argument rather than presuming a continuous choice of
sign. Work at finitely many test edges near fixed protected bulk
points, stopped before they lose a larger protection buffer, and
before the observables exceed the deterministic interior bound.
At a genuine exploration decision there are two possible next
pairings. Each next domain and its observable are determined by the
revealed past and the hypothetical pairing. The two conditional
probabilities are at least \(p_0>0\), by
Lemma~\ref{geo:domain-markov}.

Fix a tolerance \(\varepsilon>0\). Stop immediately before the first
decision at which either hypothetical next square changes by more
than \(\varepsilon\), or either hypothetical next domain loses the
smaller protection buffer, while the prevalue remains in the larger
one. This is a predictable decision index: no future continuation
is used to test the two hypothetical domains. Select one bad
alternative measurably if there are two. Conditional on reaching
this index, the probability of selecting a bad alternative is at
least \(p_0\). Consequently
\begin{equation}
 p_0\Prob\{\text{a bad-alternative stop occurs}\}
 \le
 \Prob\{\text{an observed bad step occurs}\}.
 \label{geo:bad-alternative}
\end{equation}
The right side tends to zero by the preceding square convergence,
the protected-domain convergence, and the fact that individual
cell steps have vanishing capacity increment on a protected finite
capacity interval. There is no sum over the number of decisions.

Suppose now that the current value is \(f\), with \(|f|\ge a>0\),
and both next squares are within \(\varepsilon\) of \(f^2\).
For \(\varepsilon\ll a^2\), each next value lies within
\(K\varepsilon/a\) of either \(f\) or \(-f\). If one alternative
were near \(-f\), the real projection of the conditional average
onto the direction \(f\) would be at most
\[
 |f|-2p_0|f|+K\varepsilon/a.
 \]
For \(\varepsilon\ll p_0a^2\), this contradicts the martingale
identity. Thus neither alternative changes sheets where the square
is bounded away from zero.

This statement gives tightness through zeros as well. On a short
time interval, uniform continuity of the limiting square implies
one of two possibilities. Either all its amplitudes are small,
in which case every square root is small, or its values stay in
a small disc not containing zero, in which case the signed field
stays on one continuous root branch and has a small oscillation.
Rapidly moving zeros cannot produce a fixed-size oscillation:
such an oscillation would require either a nonzero sheet switch
or a visible oscillation of the square. Deterministic forced
steps add no martingale decision and are handled by the same
square convergence. Combining this argument with
\eqref{geo:bad-alternative} proves continuity tightness at each
protected test point.

Take countably many test points and rational protected
neighborhoods, together with a countable collection of stopping
levels. Extract jointly with the curve and driver. Spatial
compactness of the fermions up to one common sign, furnished by
Proposition~\ref{prop:inputs}, identifies the limits on overlapping
neighborhoods and gives the whole field. The spatial equation
passes to the limit. Its interior elliptic estimates give all
spatial derivatives on smaller compact sets. In particular the
field is locally continuous in space and time, with the asserted
square.

Finally, before any tighter protected stop, the discrete fields
are bounded martingales: their conditional-expectation
interpretation is valid at every still active test edge and their
terminal phase is unchanged. Multiply each martingale identity
by a bounded continuous function of finitely many earlier
curve--field evaluations and pass to the joint limit. The uniform
bound gives uniform integrability. A monotone-class argument
gives the conditional expectation identity in the joint natural
filtration. Stopping at generic levels of continuous protection
quantities, or first using a tighter buffer and then exhausting,
removes any discontinuity of the stopping operation.
This proves precisely the stated local martingale assertion on
the live domain. It gives no assertion about a zero extension
through swallowing, and none will be used.
\end{proof}

\begin{remark}[Measurability and the use of the filtration]
\label{geo:filtration}
The scalar field \(h_t\) and its square derivative are functions of
the current marked domain. Their interior data are kernel-continuous,
and integrals of nonnegative functions of them are measurable by
compact exhaustion. A signed lift need not be selected as a
continuous deterministic function of every possible curve; the
joint-limit construction in Proposition~\ref{geo:martingale} is
enough for stochastic calculus. The geometric estimates in
Corollary~\ref{geo:restart} are needed only in the natural
driver filtration. Their unconditional integrability consequences
remain valid on adjoining the signed field, without an assertion
that arbitrary filtration enlargements preserve conditional
crossing estimates.
\end{remark}
